\begin{figure}[t]
\centering

% --------------------------------------------------
% Shared styles
% --------------------------------------------------
\tikzset{
    nbhfig/.style={
        >=Latex,
        line cap=round,
        line join=round,
        every node/.style={font=\small},
        lattice/.style={circle, fill=black!30, inner sep=1.05pt},
        feasiblept/.style={circle, fill=black, inner sep=1.45pt},
        lpregion/.style={draw=none, fill=blue!8},
        constraint/.style={blue!70!black, line width=1.0pt},
        solLP/.style={circle, fill=blue!80!black, inner sep=2pt},
        incpt/.style={circle, fill=red!70!black, inner sep=2pt},
        rinspt/.style={circle, fill=orange!80!black, inner sep=2pt},
        vnspt/.style={circle, fill=green!55!black, inner sep=2pt},
        shakept/.style={circle, fill=purple!80!black, inner sep=2pt},
        rinsband/.style={fill=orange!18, draw=orange!70!black, line width=0.8pt},
        nbhone/.style={draw=purple!70!black, dashed, line width=0.9pt},
        nbhtwo/.style={draw=green!60!black, dashed, line width=0.9pt}
    }
}

% --------------------------------------------------
% Shared zoom window
% --------------------------------------------------
\def\zxmin{2.45}
\def\zxmax{5.15}
\def\zymin{-0.0}
\def\zymax{2.55}

% --------------------------------------------------
% Shared geometry
% --------------------------------------------------
\newcommand{\DefineNeighborhoodGeometry}{%
    \def\eps{0.20}%
    \pgfmathsetmacro{\aone}{(-1 - 2*\eps)/4}%
    \pgfmathsetmacro{\bone}{3 + \eps}%
    \pgfmathsetmacro{\atwo}{(-3 + 2*\eps)/2}%
    \pgfmathsetmacro{\btwo}{-\atwo*(5+\eps)}%
    \pgfmathsetmacro{\xLP}{(\btwo-\bone)/(\aone-\atwo)}%
    \pgfmathsetmacro{\yLP}{\aone*\xLP+\bone}%
    \pgfmathsetmacro{\xD}{-\btwo/\atwo}%
    \pgfmathsetmacro{\rinsAymax}{3.87*\atwo+\btwo}%
    \pgfmathsetmacro{\rinsBymax}{4.13*\atwo+\btwo}%


    % Points used in the illustration
    \coordinate (xbar)  at (\xLP,\yLP);   % fractional LP solution
    \coordinate (xinc)  at (4,0);         % incumbent
    \coordinate (xrins) at (4,1);         % RINS improved solution
    \coordinate (xshake) at (4,1);        % VNS shaking point
    \coordinate (xvns)  at (3,2);         % VNS improved solution
    \coordinate (rinsmaxA) at (3.87,\rinsAymax);
    \coordinate (rinsmaxB) at (4.13,\rinsBymax);
}

% --------------------------------------------------
% Draw feasible region and lattice
% --------------------------------------------------
\newcommand{\DrawNeighborhoodBase}{%
    % LP relaxation
    \filldraw[lpregion]
        (0,0) --
        (\xD,0) --
        (\xLP,\yLP) --
        (0,\bone) --
        cycle;

    % Constraint boundaries
    \draw[constraint] (0,\bone) -- (\xLP,\yLP);
    \draw[constraint] (\xLP,\yLP) -- (\xD,0);
    \draw[constraint] (0,0) -- (\xD,0);
    \draw[constraint] (0,0) -- (0,\bone);

    % Integer lattice points
    \foreach \x in {0,...,5}{%
      \foreach \y in {0,...,4}{%
        \pgfmathtruncatemacro{\ok}{%
          ifthenelse(%
            \y <= \aone*\x+\bone + 0.001 &&%
            \y <= \atwo*\x+\btwo + 0.001,%
            1,0)%
        }%
        \ifnum\ok=1
          \node[feasiblept] at (\x,\y) {};%
        \else
          \node[lattice] at (\x,\y) {};%
        \fi
      }%
    }%
}

% --------------------------------------------------
% Draw zoom axes
% --------------------------------------------------
\newcommand{\DrawNeighborhoodZoomAxes}{%
    \draw[->, thick]
      (\zxmin,\zymin) -- (\zxmax+0.18,\zymin) node[below right] {$x_1$};
    \draw[->, thick]
      (\zxmin,\zymin) -- (\zxmin,\zymax+0.18) node[above left] {$x_2$};

    \foreach \x/\lab in {3/3,4/4,5/5}{%
      \draw (\x,\zymin) -- (\x,\zymin-0.035) node[below] {\scriptsize \lab};%
    }%
    \foreach \y/\lab in {0/0,1/1,2/2}{%
      \draw (\zxmin,\y) -- (\zxmin-0.035,\y) node[left] {\scriptsize \lab};%
    }%
}

% ==================================================
% Global view
% ==================================================
\begin{subfigure}[t]{0.315\textwidth}
\centering
\vspace{0pt}%
\begin{tikzpicture}[nbhfig, x=0.80cm, y=0.80cm]

\DefineNeighborhoodGeometry

\path[use as bounding box] (-0.35,-0.38) rectangle (6.20,4.55);

% Axes
\draw[->, thick] (-0.25,0) -- (5.75,0) node[below right] {$x_1$};
\draw[->, thick] (0,-0.25) -- (0,4.20) node[above left] {$x_2$};

\DrawNeighborhoodBase

% Zoom window
\draw[black!45, dashed, rounded corners=1pt]
    (\zxmin,\zymin) rectangle (\zxmax,\zymax);

% Main points
\node[solLP] at (xbar) {};
\node[blue!80!black, above left] at (xbar) {$\bar{x}$};

\node[incpt] at (xinc) {};
\node[red!70!black, above left] at (xinc) {$x^{\mathrm{inc}}$};

\node[rinspt] at (xrins) {};
\node[orange!80!black, above right, xshift=1mm] at (xrins) {$x^{\mathrm{RINS}}$};

\node[vnspt] at (xvns) {};
\node[green!45!black, above left] at (xvns) {$x^{\mathrm{VNS}}$};

\end{tikzpicture}
\end{subfigure}%
\hfill
% ==================================================
% Zoom: RINS
% ==================================================
\begin{subfigure}[t]{0.315\textwidth}
\centering
\vspace{0pt}%
\begin{tikzpicture}[nbhfig, x=1.55cm, y=1.25cm]

\DefineNeighborhoodGeometry

\path[use as bounding box]
  (\zxmin-0.22,\zymin-0.25) rectangle (\zxmax+0.38,\zymax+0.35);

\DrawNeighborhoodZoomAxes

\begin{scope}
  \clip (\zxmin,\zymin) rectangle (\zxmax,\zymax);

  \DrawNeighborhoodBase

  % RINS neighborhood: fix x1 = 4
  %\filldraw[rinsband]
  %  (3.87,\zymin) rectangle (4.13,\zymax);
  \filldraw[rinsband]
     (3.87,\zymin) --
     (4.13,\zymin) --
     (rinsmaxB) --
     (rinsmaxA) --
     cycle;


  % Points
  \node[solLP] at (xbar) {};
  \node[blue!80!black, above left] at (xbar) {$\bar{x}$};

  \node[incpt] at (xinc) {};
  \node[red!70!black, above right] at (xinc) {$x^{\mathrm{inc}}$};

  \node[rinspt] at (xrins) {};
  \node[orange!80!black, left, xshift=-1mm] at (xrins) {$x^{\mathrm{RINS}}$};

  % Guides / arrows
  \draw[orange!80!black, dashed, ->] (xinc) -- (xrins);
\end{scope}

\node[orange!80!black, align=center, xshift=2.5mm] at (4.45,2.20)
  {\scriptsize fix $x_1=4$};

\end{tikzpicture}
\end{subfigure}%
\hfill
% ==================================================
% Zoom: VNS
% ==================================================
\begin{subfigure}[t]{0.315\textwidth}
\centering
\vspace{0pt}%
\begin{tikzpicture}[nbhfig, x=1.55cm, y=1.25cm]

\DefineNeighborhoodGeometry

\path[use as bounding box]
  (\zxmin-0.22,\zymin-0.25) rectangle (\zxmax+0.38,\zymax+0.35);

\DrawNeighborhoodZoomAxes

\begin{scope}
  \clip (\zxmin,\zymin) rectangle (\zxmax,\zymax);

  \DrawNeighborhoodBase

  % Expanding neighborhoods around incumbent
  \draw[nbhone] (xinc) ellipse [x radius=1, y radius=1];
  

  % Points
  \node[incpt] at (xinc) {};
  \node[red!70!black, above right] at (xinc) {$x^{\mathrm{inc}}$};

  \node[shakept] at (xshake) {};
  \node[purple!80!black, above left] at (xshake) {$x'$};

  \draw[nbhtwo] (xshake) ellipse [x radius={sqrt(2)}, y radius={sqrt(2)})];

  \node[vnspt] at (xvns) {};
  \node[green!45!black, above] at (xvns) {$x^{\mathrm{VNS}}$};

  % Path
  \draw[purple!80!black, dashed, ->] (xinc) -- (xshake);
  \draw[green!55!black, dashed, ->] (xshake) -- (xvns);
\end{scope}

\node[purple!80!black] at (4.92,0.92) {\scriptsize $\mathcal{N}_1(x^{\mathrm{inc}})$};
\node[green!55!black] at (4.95,1.72) {\scriptsize $\mathcal{N}_2(x')$};

\end{tikzpicture}
\end{subfigure}

\caption{Illustration of neighborhood-based improvement heuristics. RINS defines a restricted subproblem by fixing variables on which the incumbent and the rounded LP solution agree, whereas VNS explores progressively larger neighborhoods around the incumbent through shaking and local improvement.}
\label{fig:neighborhood_search_rins_vns}
\end{figure}